\documentclass[a4,11pt]{aleph-notas}
% Se puede ver la documentación aquí:
% https://github.com/alephsub0/LaTeX_aleph-notas

% -- Paquetes adicionales
\usepackage{enumitem}
\usepackage{url}
\usepackage{array}
\usepackage{booktabs}
\usepackage{longtable}
\hypersetup{
    urlcolor=blue,
    linkcolor=blue,
}

% -- Datos
\institucion{Facultad de Ciencias Exactas, Naturales y Ambientales}
\carrera{Catálogo STEM}
\asignatura{Álgebra lineal}
\tema{Clase 04: Determinantes}
\autor{Andrés Merino}
\fecha{Periodo 2026-01}

\logouno[0.16\textwidth]{Logos/logoPUCE_04_ac}
\definecolor{colortext}{HTML}{1957d1}
\definecolor{colordef}{HTML}{1957d1}
\fuente{montserrat}

\begin{document}

\encabezado

%%%%%%%%%%%%%%%%%%%%%%%%%%%%%%%%%%%%%%%%
\section*{Resultado de Aprendizaje}
%%%%%%%%%%%%%%%%%%%%%%%%%%%%%%%%%%%%%%%%

%%%%%%%%%%%%%%%%%%%%%%%%%%%%%%%%%%%%%%%%
\subsection*{RdA de la asignatura:}
\begin{itemize}[leftmargin=*]
    \item \textbf{RdA 1:} Comprender los conceptos fundamentales del Álgebra Lineal, incluyendo el estudio de matrices, determinantes, sistemas de ecuaciones lineales y espacios vectoriales, destacando su importancia en el análisis y resolución de problemas matemáticos.
    \item \textbf{RdA 2:} Resolver problemas prácticos y teóricos relacionados con el Álgebra Lineal, utilizando técnicas de cálculo con matrices, determinantes y sistemas de ecuaciones, así como en el análisis de espacios vectoriales.
\end{itemize}


%%%%%%%%%%%%%%%%%%%%%%%%%%%%%%%%%%%%%%%%
\subsection*{RdA de la actividad:}
\begin{itemize}[leftmargin=*]
    \item Comprender el significado algebraico y geométrico del determinante.
    \item Aplicar propiedades del determinante para simplificar cálculos en matrices cuadradas.
    \item Identificar condiciones de singularidad y su relación con la invertibilidad.
    \item Verificar resultados analíticos con apoyo de herramientas computacionales.
\end{itemize}

%%%%%%%%%%%%%%%%%%%%%%%%%%%%%%%%%%%%%%%%
\section*{Introducción}
%%%%%%%%%%%%%%%%%%%%%%%%%%%%%%%%%%%%%%%%

%%%%%%%%%%%%%%%%%%%%%%%%%%%%%%%%%%%%%%%%
\paragraph{Control de avance y clase invertida:}
Antes de iniciar la clase, se aplicará el control de avance, tomando un ejercicio de multiplicación de matrices y otro de cálculo de determinantes.

\paragraph{Pregunta inicial:}
¿Qué información guarda el determinante de una matriz? ¿Cómo se relaciona con la geometría y la invertibilidad?

%%%%%%%%%%%%%%%%%%%%%%%%%%%%%%%%%%%%%%%%
\section*{Desarrollo}
%%%%%%%%%%%%%%%%%%%%%%%%%%%%%%%%%%%%%%%%

%%%%%%%%%%%%%%%%%%%%%%%%%%%%%%%%%%%%%%%%
\subsection*{Actividad 1: Determinantes y sus propiedades}

La clase inicia con una visión conjunta del significado del determinante en distintos tamaños de matrices. Luego, mediante clase magistral y práctica guiada, se formalizan propiedades algebraicas y su interpretación para analizar invertibilidad y resolver problemas.

\paragraph{¿Cómo lo haremos?}
\begin{itemize}[leftmargin=*]
    \item \textbf{Clase invertida -- Visión conjunta:} se repasa el cálculo de determinantes para matrices de orden 1, 2 y 3.
    \item \textbf{Clase magistral:} se explican las propiedades del determinante. Se usa el resumen \href{https://fcena-puce.github.io/AlgebraLineal-05-N0048/01Resumen/Resumen04.pdf}{Resumen04.pdf}.
    \item \textbf{Resolución de ejercicios:} se guiará a los estudiantes en la solución de ejercicios del documento \href{https://fcena-puce.github.io/AlgebraLineal-05-N0048/04Ejercicios/Ejercicios04.pdf}{Ejercicios04.pdf}, aplicando las propiedades vistas.
    \item \textbf{Visualización de video:} se recomendará el estudio con el video: \href{https://youtu.be/Ip3X9LOh2dk?si=Z8KpgrPdLpy8l48A}{El determinante | Capítulo 6, Esencia del álgebra lineal.}.
\end{itemize}

\paragraph{Verificación de aprendizaje:}
\begin{enumerate}[leftmargin=*]
    \item ¿Qué propiedades tiene el determinante?
    \item ¿Cuándo el determinante de una matriz es cero? ¿Qué implica eso?
    \item ¿Cómo cambia el determinante si permutamos dos filas?
\end{enumerate}

%%%%%%%%%%%%%%%%%%%%%%%%%%%%%%%%%%%%%%%%
\section*{Cierre}
%%%%%%%%%%%%%%%%%%%%%%%%%%%%%%%%%%%%%%%%

\paragraph{Tarea:}
Resolver del libro \href{https://catalogobiblioteca.puce.edu.ec/cgi-bin/koha/opac-detail.pl?biblionumber=86081}{Fundamentos de álgebra lineal de Larson}, sección 3.1, los ejercicios: 1, 5, 9, 11, 27, 49, 51; sección 3.2, los ejercicios: 1, 3, 5, 7, 12, 13; sección 3.3, los ejercicios: 53, 56.

\paragraph{Pregunta de investigación:}
\begin{enumerate}[leftmargin=*]
    \item ¿Qué es la regla de Cramer y cómo se relaciona con el determinante?
    \item ¿Cómo se puede usar el determinante para encontrar el área de un paralelogramo definido por dos vectores en $\R^2$?
    \item ¿Cómo se puede usar el determinante para encontrar el volumen de un paralelepípedo definido por tres vectores en $\R^3$?
\end{enumerate}

\paragraph{Para la próxima clase:}  
Leer la sección 4.1 del libro \href{https://catalogobiblioteca.puce.edu.ec/cgi-bin/koha/opac-detail.pl?biblionumber=86081}{Fundamentos de álgebra lineal de Larson}.




\end{document}
